\chapter{Marco Teórico}

\section{Microscopía y caracterización a escala micrométrica}

La microscopía constituye un conjunto de técnicas orientadas a la observación y caracterización de estructuras cuyo tamaño se encuentra por debajo del límite de resolución del ojo humano. En el contexto de la ingeniería y la investigación científica, la microscopía no se limita a la visualización cualitativa, sino que se emplea como herramienta para la medición, comparación y análisis morfológico de estructuras a escala micrométrica y submicrométrica.

En aplicaciones biológicas y biomiméticas, la microscopía permite identificar características estructurales tales como tamaños celulares, bordes, distribuciones espaciales y relaciones geométricas básicas. Estas observaciones constituyen una base fundamental para estudios exploratorios, educativos y de investigación básica, sin implicar necesariamente un uso clínico o diagnóstico.

Desde un punto de vista instrumental, los sistemas de microscopía pueden clasificarse según su arquitectura óptica, el tipo de detector empleado y el grado de procesamiento computacional involucrado en la formación de la imagen. En este marco, los sistemas tradicionales basados en lentes conviven con enfoques alternativos que priorizan la simplicidad del hardware y delegan parte de la interpretación visual al procesamiento posterior.

\section{Fundamentos físicos de la microscopía sin lentes}

La microscopía sin lentes (\textit{lensless microscopy}) es un enfoque en el cual la formación de imagen no depende de sistemas ópticos refractivos tradicionales, sino de la interacción directa entre la luz y la muestra, registrada por un detector fotosensible ubicado a corta distancia.

En configuraciones de microscopía de contacto, la muestra se coloca a una distancia del orden de micrómetros del plano del sensor, de modo que la imagen resultante está dominada por fenómenos de absorción, dispersión cercana y difracción. En este régimen, la luz transmitida o dispersada por la muestra proyecta patrones de intensidad directamente sobre el sensor, los cuales contienen información estructural de la muestra.

La ausencia de lentes reduce significativamente la complejidad mecánica y óptica del sistema, permitiendo arquitecturas compactas, robustas y de bajo consumo energético. Sin embargo, esta simplificación introduce limitaciones inherentes relacionadas con la resolución efectiva, la sensibilidad al acoplamiento óptico muestra–sensor y la influencia de efectos difractivos.

\section{Sensor CMOS como sistema de adquisición}

Los sensores CMOS (Complementary Metal-Oxide-Semiconductor) se han consolidado como una tecnología madura y ampliamente disponible para la adquisición de imágenes digitales. Su bajo consumo de potencia, alta integración y compatibilidad con sistemas embebidos los convierten en una opción adecuada para plataformas con recursos limitados.

En el contexto de microscopía sin lentes, el sensor CMOS actúa como el plano de detección donde se registra la distribución espacial de intensidad luminosa generada por la interacción de la luz con la muestra. Cada píxel del sensor integra fotones incidentes durante un intervalo de tiempo definido, produciendo una señal digital proporcional a la intensidad local.

El tamaño físico del píxel, la eficiencia cuántica, el ruido electrónico y la resolución nativa del sensor establecen límites fundamentales a la capacidad del sistema para resolver estructuras pequeñas.

\begin{table}[H]
\centering
\caption{Características técnicas relevantes del sensor CMOS utilizado}
\begin{tabular}{|l|l|}
\hline
\textbf{Parámetro} & \textbf{Valor típico} \\
\hline
Tipo de sensor & CMOS (Commercial Off-The-Shelf) \\
Resolución nativa & 2592 × 1944 píxeles \\
Tamaño de píxel & $\sim 1.4\,\mu m$ \\
Área activa & $\sim 3.7\,mm \times 2.7\,mm$ \\
Profundidad de bits & 10–12 bits \\
Control de exposición & Manual \\
Control de ganancia & Manual \\
Filtro de color & Matriz Bayer (RGB) \\
Interfaz de datos & CSI / interfaz digital directa \\
\hline
\end{tabular}
\end{table}

\section{Resolución, difracción y límites físicos del sistema}

La resolución espacial de un sistema de imagen se define como la capacidad de distinguir dos estructuras cercanas como entidades separadas. En microscopía sin lentes, la resolución efectiva está influenciada por:

\begin{itemize}
    \item Tamaño físico del píxel del sensor.
    \item Distancia muestra–sensor.
    \item Longitud de onda de la iluminación.
    \item Respuesta espacial del detector.
\end{itemize}

Cuando las dimensiones de la muestra son comparables al límite de resolución del sistema, las estructuras no se observan como objetos bien definidos, sino como patrones difusos dominados por fenómenos de difracción. Este comportamiento es esperable y responde a principios físicos fundamentales de propagación de la luz.

\section{Escala espacial y calibración dimensional}

La obtención de mediciones dimensionales a partir de imágenes digitales requiere establecer una relación confiable entre el espacio físico y el espacio discreto representado por los píxeles del sensor. Esta relación, comúnmente expresada como escala espacial en unidades de micrómetros por píxel ($\mu m/pixel$), debe validarse experimentalmente.

En desarrollos experimentales basados en sensores CMOS, la escala inicial puede estimarse a partir de parámetros nominales del sensor. No obstante, esta estimación debe contrastarse mediante procedimientos independientes antes de utilizar el sistema para caracterización dimensional de muestras reales.

La calibración dimensional constituye, por tanto, un paso crítico para asegurar la trazabilidad y coherencia metrológica del sistema.

\section{Objetos patrón micrométricos para validación}

Los objetos patrón micrométricos se emplean como referencias físicas conocidas para validar la escala espacial y la capacidad de medición de un sistema de imagen. Entre los más utilizados se encuentran microesferas poliméricas y patrones litográficos con geometrías definidas.

Las microesferas de poliestireno presentan diámetros nominales certificados y una dispersión de tamaños conocida, lo que permite evaluar la coherencia dimensional del sistema.

\begin{table}[H]
\centering
\caption{Características típicas de microesferas de poliestireno}
\begin{tabular}{|l|l|}
\hline
\textbf{Parámetro} & \textbf{Valor típico} \\
\hline
Material & Poliestireno \\
Diámetro nominal & 2 $\mu m$ \\
Coeficiente de variación & $\sim 5\%$ \\
Forma & Esférica \\
Medio de suspensión & Solución acuosa \\
Índice de refracción & $\sim 1.59$ \\
\hline
\end{tabular}
\end{table}

\begin{table}[H]
\centering
\caption{Comparación cualitativa de objetos patrón}
\begin{tabular}{|l|c|c|}
\hline
\textbf{Característica} & \textbf{Microesferas} & \textbf{Patrones litográficos} \\
\hline
Geometría & Esférica & Lineal / reticular \\
Dimensión conocida & Nominal (distribuida) & Certificada \\
Influencia de difracción & Alta cerca del límite & Menor \\
Robustez metrológica & Media & Alta \\
Adecuación para $\mu m/pixel$ & Complementaria & Principal \\
\hline
\end{tabular}
\end{table}

\section{Microscopía computacional como apoyo a la interpretación}

La microscopía computacional engloba técnicas orientadas a mejorar la interpretabilidad de imágenes mediante procesamiento posterior. Estas técnicas no sustituyen la adquisición física de la imagen, sino que actúan como herramientas complementarias para realce visual y análisis exploratorio.

En el presente trabajo, el procesamiento computacional se emplea como apoyo a la interpretación de estructuras ya presentes en la imagen base, reconociendo explícitamente que no incrementa la resolución óptica real del sistema ni genera información física adicional.

Este enfoque garantiza la trazabilidad dimensional y mantiene una separación clara entre adquisición física y postprocesamiento digital.
